\begin{abstract}

    % Francais
    Ce travail de Bachelor porte sur la conception complète d'un lecteur audio portable, du schéma électrique au routage du circuit imprimé, puis au firmware embarqué et aux mesures de validation. L'objectif était de réaliser un appareil compact, autonome et dédié à la seule lecture de musique, réellement fabriqué et fonctionnel. La démarche suit un déroulement descendant. Une analyse fonctionnelle et des spécifications techniques ont d'abord été établies, puis la carte quatre couches a été conçue sous Altium Designer autour d'un microcontrôleur ESP32, le firmware temps réel a été développé sous ESP-IDF avec FreeRTOS et le boîtier a été imprimé en 3D. L'appareil terminé a ensuite été caractérisé par une campagne de mesure. Il lit des fichiers MP3 et WAV depuis une carte microSD jusqu'à 24 bits et 192\,kHz, restitue le son sur un jack 3,5\,mm ou en Bluetooth A2DP, s'utilise via un écran OLED et des menus, et se charge et se programme par un connecteur USB-C unique. L'interface graphique, rendue par un moteur d'affichage développé sur mesure, offre une barre de statut permanente, la liste des pistes, des menus de réglages et des pages de statistiques. Dix des onze exigences minimales du cahier des charges sont pleinement satisfaites. La sortie filaire présente un THD+N de $-94{,}7$\,dB et un SNR de 106,2\,dB, et la liaison Bluetooth reste stable sur une vingtaine de mètres. L'autonomie atteint 9\,h\,14 en lecture filaire et 5\,h\,48 en Bluetooth, sous l'objectif de dix heures. Le prototype valide l'architecture retenue et constitue la base solide d'une révision future dont les améliorations sont identifiées.
    
    %% L'asterisme est un signe typographique en forme d'étoile, utilisé pour marquer une pause dans un texte ou pour séparer des paragraphes. Il est souvent utilisé pour indiquer un changement de scène dans un récit. Bien qu'il se fasse rare dans la typographie moderne, c'est un symbole de choix pour séparer les différentes langues du résumé de thèse.
    \asterism
    
    % English
    This Bachelor's thesis covers the complete design of a portable audio player, from the electrical schematic to the printed circuit board layout, then to the embedded firmware and the validation measurements. The goal was to build a compact, self-contained device dedicated solely to music playback, actually manufactured and functional. The work follows a top-down approach. A functional analysis and technical specifications were established first, then the four-layer board was designed in Altium Designer around an ESP32 microcontroller, the real-time firmware was developed with ESP-IDF and FreeRTOS and the enclosure was 3D printed. The finished device was then characterized through a measurement campaign. It plays MP3 and WAV files from a microSD card up to 24 bits and 192\,kHz, outputs audio through a 3.5\,mm jack or over Bluetooth A2DP, is operated through an OLED display and a menu system, and is both charged and programmed through a single USB-C connector. The graphical interface, drawn by a custom rendering engine, provides a permanent status bar, a track list, settings menus and statistics pages. Ten of the eleven minimum requirements are fully met. The wired output achieves a THD+N of $-94.7$\,dB and an SNR of 106.2\,dB, and the Bluetooth link remains stable over some twenty meters. Battery life reaches 9\,h\,14 in wired playback and 5\,h\,48 over Bluetooth, below the ten-hour target. The prototype validates the chosen architecture and provides the basis for a revision whose improvements are identified.
    
\end{abstract}